% R008 — Guide to Cultivating Complex Analysis (Bahasa Indonesia)
% Reader-facing translation unit: source/ca-v1.9/ca.tex lines 1054–1072.
% Source release: v1.9 (commit a4d25d9ba37a486a5a020219eb8bd4e6602fd14c).
% Translation provenance: OpenAI Codex gpt-5.6-sol, Ultra, acting on the user's request.
% Preserve source footnote formula, label, notation, and matrix structure.

\subsection{Representasi matriks bilangan kompleks}

Karena $\C$ adalah $\R^2$, kita dapat memandang
$\C$ sebagai ruang vektor real berdimensi dua dengan mengabaikan perkalian
kompleks.
Basis standarnya adalah $1$ dan $i$.
Untuk memasukkan kembali perkalian ke dalam pembahasan, kita memandang
operator linear pada $\R^2$.
Diberikan bilangan kompleks $\xi$, pemetaan $z \mapsto \xi z$ merupakan
operator linear-real\footnote{$L \colon \C \to \C$ bersifat linear-real jika
$L(a z + b w) = aL(z)+bL(z)$ untuk semua $a,b \in \R$ dan $z,w \in \C$.}.
Suatu operator linear-real pada $\R^2$ diberikan oleh matriks real
berukuran $2 \times 2$.
Tepatnya, bilangan kompleks $\xi = a+ib$ dapat direpresentasikan oleh matriks
real berukuran $2 \times 2$
\begin{equation} \label{eq:complexnumbermatrix}
\begin{bmatrix}
a & -b \\
b & a
\end{bmatrix} .
\end{equation}
